\chapter{Complexity — Unary Logspace}
%%% generated by the blueprint pipeline from TCSlib.Complexity.SpaceComplexity.UnaryLogspace
\section{Overview}

Declarations of the Lean module \texttt{TCSlib.Complexity.SpaceComplexity.UnaryLogspace}.

\section{Declarations}

\begin{definition}[Unary-indexed bit language]
\lean{Complexity.uBit}
\label{Complexity.uBit}
\leanok
\uses{Turing.pairEncode}
For a function $g$ assigning a Boolean list to each natural number, define a language
over the Boolean alphabet consisting of precisely those pair-encoded words
$(1^n,	ext{bits}(i))$ for which the $i$th entry of $g(n)$ is true, taking an
out-of-range entry to be false.
\end{definition}

\begin{definition}[Unary-indexed length language]
\lean{Complexity.uLen}
\label{Complexity.uLen}
\leanok
\uses{Turing.pairEncode}
For a function $g$ assigning a Boolean list to each natural number, define a language
over the Boolean alphabet consisting of precisely those pair-encoded words
$(1^n,	ext{bits}(i))$ for which $i$ is less than the length of $g(n)$.
\end{definition}

\begin{definition}[Unary logspace]
\lean{Complexity.UnaryLogspace}
\label{Complexity.UnaryLogspace}
\leanok
\uses{Complexity.LOGSPACE, Complexity.uBit, Complexity.uLen}
A function $g$ from natural numbers to Boolean lists is unary logspace if both its
unary-indexed bit language and its unary-indexed length language belong to
$	extsf{LOGSPACE}$.
\end{definition}

\begin{lemma}[Membership in the encoded index set]
\lean{Complexity.mem_unaryIdx}
\label{Complexity.mem_unaryIdx}
\leanok
\uses{Complexity.LogProg.bits_injective, Turing.pairEncode, Turing.pairEncode_replicate_inj}
\difficulty{4}
Let $p$ be a predicate on pairs of natural numbers, and let $n,i\in\mathbb{N}$. The
string formed by pairing $n$ copies of the bit $1$ with the binary bit string of $i$
belongs to the set of all strings that encode some pair $(r,s)$ satisfying $p(r,s)$ if
and only if $p(n,i)$ holds.
\end{lemma}

\begin{lemma}[Membership test for a unary-indexed bit language]
\lean{Complexity.mem_uBit}
\label{Complexity.mem_uBit}
\leanok
\uses{Complexity.mem_unaryIdx, Complexity.uBit, Turing.pairEncode}
\difficulty{1}
For every function $g$ from natural numbers to binary strings and all natural numbers
$n,i$, the self-delimiting encoding of $n$ copies of $1$ paired with the binary
expansion of $i$ belongs to the language determined by $g$ exactly when the bit of
$g(n)$ at position $i$ is $1$, with positions counted from $0$ and out-of-range
positions treated as $0$.
\end{lemma}

\begin{lemma}[Membership test for a unary-indexed length language]
\lean{Complexity.mem_uLen}
\label{Complexity.mem_uLen}
\leanok
\uses{Complexity.mem_unaryIdx, Complexity.uLen, Turing.pairEncode}
\difficulty{1}
For every function $g$ from natural numbers to binary strings and all natural numbers
$n,i$, the self-delimiting encoding of $n$ copies of $1$ paired with the binary
expansion of $i$ belongs to the language determined by the lengths of the values of $g$
exactly when $i<|g(n)|$.
\end{lemma}

\begin{lemma}[Unary-index encodings are valid]
\lean{Complexity.validPlain_of_mem_unaryIdx}
\label{Complexity.validPlain_of_mem_unaryIdx}
\leanok
\uses{Complexity.LogProg.ValidPlain, Complexity.LogProg.canon_bits, Turing.pairEncode}
\difficulty{4}
Let $p$ be a predicate on pairs of natural numbers, and let $y$ be a bit string. If
there are natural numbers $n$ and $i$ such that $y$ is the pairing of the length-$n$
all-ones string with the binary expansion of $i$, and $p(n,i)$ holds, then $y$ is a
valid plain encoding.
\end{lemma}

\begin{definition}[Unary extension of a list function]
\lean{Complexity.unaryExt}
\label{Complexity.unaryExt}
\leanok
Given a function $g$ from natural numbers to Boolean lists, define its extension to
Boolean lists by sending a list $x$ to $g(|x|)$ when every entry of $x$ is true, and to
the empty list otherwise.
\end{definition}

\begin{lemma}[Unary extension on a unary word]
\lean{Complexity.unaryExt_replicate}
\label{Complexity.unaryExt_replicate}
\leanok
\uses{Complexity.unaryExt}
\difficulty{1}
For every function $g:\mathbb N\to\{0,1	ext{-strings}	ext{, including the empty
string}
ightarrow$ and every $n\in\mathbb N$, the extension of $g$ evaluated at the word
of $n$ ones equals $g(n)$.
\end{lemma}

\begin{lemma}[Unary extension of an indexed language]
\lean{Complexity.indexLang_unaryExt}
\label{Complexity.indexLang_unaryExt}
\leanok
\uses{Complexity.indexLang, Turing.pairEncode}
\difficulty{4}
Let $p$ be a predicate on bit strings and natural numbers. Suppose that for every bit
string $x$ and natural number $i$, $p(x,i)$ holds exactly when $x$ is the all-ones
string of length $|x|$ and $p(x,i)$ holds. Let $q$ be a predicate on natural numbers,
and suppose that for every natural numbers $n,i$, $p(1^n,i)$ holds exactly when $q(n,i)$
holds. Then the language of encodings $	ext{pairEncode}(x,	ext{bits}(i))$ satisfying
$p(x,i)$ is exactly the language of encodings $	ext{pairEncode}(1^n,	ext{bits}(i))$
satisfying $q(n,i)$.
\end{lemma}

\begin{theorem}[Unary logspace yields implicit logspace computation]
\lean{Complexity.UnaryLogspace.implicitlyLogspaceComputable}
\label{Complexity.UnaryLogspace.implicitlyLogspaceComputable}
\leanok
\uses{Complexity.ImplicitlyLogspaceComputable, Complexity.UnaryLogspace, Complexity.indexLang_unaryExt, Complexity.unaryExt}
\difficulty{4}
Let $g:\mathbb{N}\to\{0,1\}^*$ have both its bit-membership language and its
output-length language decidable in logarithmic space. Suppose there are natural numbers
$C,c$ such that $|g(n)|\leq C(n+1)^c$ for every $n$. Extend $g$ to bit strings by
returning $g(n)$ on the unary string $1^n$ and the empty string on every other input;
this extension is implicitly logspace computable, meaning its output length has a
polynomial bound and its bit-membership and length-index languages are decidable in
logarithmic space.
\end{theorem}

\begin{lemma}[Unary first-argument oracle call]
\lean{Complexity.LogProg.astep_call₁}
\leanok
\uses{Complexity.LogProg.ARM, Complexity.LogProg.astep, Complexity.LogProg.vword_unary₁, Turing.pairEncode}
\difficulty{3}
Let $A$ be a register program, $o$ an oracle function, and $l,l_1,l_0$ control states.
If the instruction at $l$ calls oracle $j$ in unary-first mode with the single argument
register $r$ and branches to $l_1$ on true and $l_0$ on false, then on input encoding a
unary word of length $n$ followed by a word $w$, one abstract step from state $l$,
register assignment $v$, and no stored answer goes to $l_1$ exactly when $o(j,\langle
1^n,\mathrm{bits}(v(r))\rangle)$ is true, and otherwise goes to $l_0$, leaving the
registers and stored answer unchanged.
\end{lemma}

\begin{definition}[Unary comparison language]
\lean{Complexity.ltLang}
\label{Complexity.ltLang}
\leanok
\uses{Turing.pairEncode}
Let $	ext{pairEncode}$ denote the pairing of two Boolean words into one Boolean word.
Define a language over the Boolean alphabet by including exactly those words that can be
written as $	ext{pairEncode}(1^n,	ext{bits}(i))$ for natural numbers $n$ and $i$ with
$i<n$, where $1^n$ is the length-$n$ word consisting entirely of true symbols and
$	ext{bits}(i)$ is the binary representation of $i$.
\end{definition}

\begin{definition}[Comparison machine labels]
\lean{Complexity.LtM.Lb}
\label{Complexity.LtM.Lb}
\leanok
Introduce eight labels for a comparison machine: start, test, compare, three
intermediate labels, yes, and no.
\end{definition}

\begin{definition}[Unary less-than machine]
\lean{Complexity.LtM.A}
\label{Complexity.LtM.A}
\leanok
\uses{Complexity.LogProg.ARM, Complexity.LtM.Lb}
Consider a two-register machine with one available decider and the eight labels start,
test, compare, three intermediate labels, yes, and no. At start it validates the input
as a unary-prefixed canonical binary word, rejecting invalid inputs; at test it calls
the available decider on the unary prefix with register 1 as its argument, branching to
no if the call accepts and to compare otherwise. At compare it accepts if the input's
binary payload equals register 1; otherwise it increments register 1, then increments
register 2 twice, and returns to test. The yes label accepts and the no label rejects.
\end{definition}

\begin{definition}[Two-register value assignment]
\lean{Complexity.LtM.vv}
\label{Complexity.LtM.vv}
\leanok
For natural numbers $a$ and $b$, define a function on the two-element index set that
assigns $a$ to index $0$ and $b$ to index $1$.
\end{definition}

\begin{lemma}[Value at the first index]
\lean{Complexity.LtM.vv_0}
\label{Complexity.LtM.vv_0}
\leanok
\uses{Complexity.LtM.vv}
\difficulty{1}
For natural numbers $a$ and $b$, define $v$ on the two indices $0$ and $1$ by $v(0)=a$
and $v(1)=b$. Then $v(0)=a$.
\end{lemma}

\begin{lemma}[Second Coordinate Value]
\lean{Complexity.LtM.vv_1}
\label{Complexity.LtM.vv_1}
\leanok
\uses{Complexity.LtM.vv}
\difficulty{1}
For natural numbers $a$ and $b$, define $v:	ext{Fin }2	obn$ by $v(0)=a$ and $v(1)=b$.
Then $v(1)=b$.
\end{lemma}

\begin{lemma}[Updating the First Coordinate]
\lean{Complexity.LtM.upd_0}
\label{Complexity.LtM.upd_0}
\leanok
\uses{Complexity.LtM.vv}
\difficulty{4}
For natural numbers $a,b,c$, let $v:	ext{Fin}(2)\to\mathbb{N}$ have values $a$ at index
$0$ and $b$ at index $1$. Updating $v$ at index $0$ to $c$ gives the function with
values $c$ at index $0$ and $b$ at index $1$.
\end{lemma}

\begin{lemma}[Updating the Second Coordinate]
\lean{Complexity.LtM.upd_1}
\label{Complexity.LtM.upd_1}
\leanok
\uses{Complexity.LtM.vv}
\difficulty{4}
For natural numbers $a,b,c$, let $v:\text{Fin}(2)\to\mathbb{N}$ have values $a$ at index
$0$ and $b$ at index $1$. Updating $v$ at index $1$ to $c$ gives the function with
values $a$ at index $0$ and $c$ at index $1$.
\end{lemma}

\begin{definition}[Indicator of the designated language]
\lean{Complexity.LtM.orc}
\label{Complexity.LtM.orc}
\leanok
\uses{Complexity.LogProg.dblLang, Turing.MultiTapeTM.indicator}
For each index in the one-element set and each finite Boolean string $V$, the function
returns the Boolean indicator of whether $V$ belongs to the designated language of
doubled strings.
\end{definition}

\begin{lemma}[Membership of encoded strings]
\lean{Complexity.LtM.orc_eq}
\label{Complexity.LtM.orc_eq}
\leanok
\uses{Complexity.LogProg.bits_injective, Complexity.LogProg.indicator_eq_decide, Complexity.LtM.orc, Turing.pairEncode, Turing.pairEncode_replicate_inj}
\difficulty{4}
For natural numbers $n$ and $v$, the Boolean oracle that tests membership in the
language of encodings $$ returns true on the pairing of the length-$n$ string of ones
with the bit string of $v$ exactly when $v=2n$.
\end{lemma}

\begin{definition}[Bounded pre-simulation configurations]
\lean{Complexity.LtM.G}
\label{Complexity.LtM.G}
\leanok
\uses{Complexity.LogProg.AConf, Complexity.LogProg.PreS, Complexity.LtM.A, Complexity.LtM.Lb}
For a Boolean string $y$ and a configuration $a$ consisting of a state, two work-tape
head positions, and an output, $G(y,a)$ holds precisely when $	ext{PreS}(A,y,a)$ holds
and each head position is at most $2(|y|+1)$.
\end{definition}

\begin{lemma}[Loop Halting and Output]
\lean{Complexity.LtM.loop}
\label{Complexity.LtM.loop}
\leanok
\uses{Complexity.LogProg.AHalt, Complexity.LogProg.AHalt.ret, Complexity.LogProg.AHalt.step, Complexity.LogProg.PreS, Complexity.LogProg.ValidPlain, Complexity.LogProg.astep, Complexity.LogProg.astep_call₁, Complexity.LogProg.bits_injective, Complexity.LogProg.canon_bits, Complexity.LogProg.dbl_replicate, Complexity.LogProg.plainWord_pairEncode, Complexity.LtM.A, Complexity.LtM.G, Complexity.LtM.Lb, Complexity.LtM.orc, Complexity.LtM.orc_eq, Complexity.LtM.upd_0, Complexity.LtM.upd_1, Complexity.LtM.vv, Complexity.LtM.vv_0, Complexity.LtM.vv_1, Turing.pairEncode, Turing.pairEncode_eq_dbl}
\difficulty{6}
Let $n$ and $i$ be natural numbers. For every pair of natural numbers $d,k$ such that
$k+d=n$ and $k\le i$, the specified machine, on input encoding the unary string of
length $n$ paired with the binary expansion of $i$, halts with output true exactly when
$i<n$. It starts in state $\mathrm{test}$ with its two registers equal to $k$ and $2k$,
respectively, and no output; at every time strictly before halting, its configuration
satisfies the stated invariant $G$.
\end{lemma}

\begin{theorem}[A comparison language is in logarithmic space]
\lean{Complexity.ltLang_mem}
\label{Complexity.ltLang_mem}
\leanok
\uses{Complexity.LOGSPACE, Complexity.LogProg.AHalt, Complexity.LogProg.AHalt.step, Complexity.LogProg.PreS, Complexity.LogProg.ValidPlain, Complexity.LogProg.arm_decides_poly, Complexity.LogProg.arun, Complexity.LogProg.astep, Complexity.LogProg.bitsVal, Complexity.LogProg.canon_bits, Complexity.LogProg.canon_eq_bits, Complexity.LogProg.dblLang, Complexity.LogProg.dblLang_mem, Complexity.LogProg.indicator_eq_decide, Complexity.LtM.A, Complexity.LtM.G, Complexity.LtM.loop, Complexity.LtM.orc, Complexity.LtM.vv, Complexity.ltLang, Complexity.mem_unaryIdx, Complexity.validPlain_of_mem_unaryIdx, Turing.MultiTapeTM.indicator, Turing.pairEncode}
\difficulty{5}
Let $L$ be the language of bit strings of the form $\langle
1^n,\operatorname{bits}(i)\rangle$ with $i<n$, where the pairing doubles each bit of the
first component, appends the separator $01$, and then appends the second component, and
$\operatorname{bits}(i)$ is the binary expansion of $i$. Then $L$ is decidable by a
deterministic finite-state multi-tape Turing machine using at most a constant multiple
of $\lfloor\log_2 m\rfloor+1$ work-tape cells on inputs of length $m$, with the value at
$m=0$ taken to be $1$.
\end{theorem}

\begin{theorem}[Unary words have logarithmic-space query languages]
\lean{Complexity.unaryLogspace_replicate}
\label{Complexity.unaryLogspace_replicate}
\leanok
\uses{Complexity.UnaryLogspace, Complexity.ltLang, Complexity.ltLang_mem, Complexity.uBit, Complexity.uLen}
\difficulty{4}
For the function $g(n)=1^n$, both languages of encoded pairs consisting of $1^n$ and the
binary representation of an index—one accepting when the indexed bit of $g(n)$ is $1$,
the other when the index is less than $|g(n)|$—are decidable in logarithmic space.
\end{theorem}
